% generated by roam.py handbook export -- не редактировать, править заметки
% 20260922144651-единственность_минимума_сильно_выпуклои_функции.org (body d3f1c9a34490)

\begin{theorem}[{Единственность минимума сильно выпуклой функции}]\label[theorem]{note:51CC880D-BF19-43E2-9599-FB31A9A7367F}
  Если $f: \R^n \rightarrow \R$~--- \hyperref[note:957FA09D-8883-40A5-B4C7-237D9F815A42]{дифференцируемая} и \hyperref[note:22E546FC-186F-4A67-85AA-2C85DCEF9021]{сильно выпуклая} функция, то у нее существует единственная точка минимума
\end{theorem}

\begin{remark}
  Дифференцируемость нужна только для существования: по \hyperref[note:1B5A3E25-56E8-4D5E-9CF6-E2731DA77604]{дифференциальному критерию} $f$ растет не медленнее параболы, поэтому минимум достигается на компактном шаре (Вейерштрасс). Единственность верна и без нее. Существование без дифференцируемости тоже верно, но доказывается сложнее~--- через непрерывность выпуклой функции на $\R^n$.
  Важно и то, что $f$ задана на всем $\R^n$: например, $f(x) = x^2$ на $(1;2)$ сильно выпукла, но минимум не достигается.
\end{remark}

\begin{proof}
  \proofpart{Существование}
  Давайте возьмем \hyperref[note:1B5A3E25-56E8-4D5E-9CF6-E2731DA77604]{дифференциальный критерий сильно выпуклой функции} для $x=0$
  \[
  f(y) \ge f(0) + \left\langle \nabla f(0); y \right\rangle + \frac{\mu}{2} \| y \|^2
  \]
  Применим \hyperref[note:D0601ACD-8C91-4A93-B0E4-AC0D529C6A5C]{КБШ} $| \left\langle x;y \right\rangle | \le \|x\| \cdot \|y\| \Leftrightarrow - \|x\| \cdot \|y\| \le   \left\langle x;y \right\rangle  \le \|x\| \cdot \|y\|$ (левую часть):
  \[
  f(y) \ge f(0) + \left\langle \nabla f(0); y\right\rangle + \frac{\mu}{2} \| y \|^2 \ge f(0) - \| \nabla f(0) \| \| y\| + \frac{\mu}{2} \| y \|^2
  \]

  Возьмем радиус $R := \frac{2 \| \nabla f(0) \|}{\mu}$, тогда, при $\|y\| > R$
  \begin{multline*}
  f(y) \ge f (0) -  \| \nabla f(0) \| \| y \| + \frac{\mu}{2} \| y \|^2 > f(0) - \| \nabla f(0) \| R + \frac{\mu}{2} R^2 = \\
   f(0) - \| \nabla f(0)\| \cdot  \frac{ 2 \| \nabla f(0)\|}{\mu} + \frac{\mu}{2} \cdot  \left( \frac{2 \| \nabla f(0) \|}{\mu} \right)^2 = \\
   f(0) - \frac{2}{\mu} \|\nabla f(0) \|^2 + \frac{2}{\mu} \| \nabla f(0) \|^2 = f(0)
  \end{multline*}
  Второе неравенство верно, так как парабола $h(r) = - \| \nabla f(0) \| r + \frac{\mu}{2} r^2$ возрастает при $r \ge \frac{\| \nabla f(0) \|}{\mu}$, а $\| y \| > R \ge \frac{\| \nabla f(0) \|}{\mu}$.
  Получили, что $f(y) > f(0)$, при $\| y \| > R$.

  Теперь рассмотрим $\| y \| \le R$~--- это замкнутый шар.
  Мы знаем, что $f$~--- дифференцируема, следовательно она непрерывна.
  Тогда по \hyperref[note:82C70263-336A-49F5-82DB-9D25C84D6A6D]{теореме Вейерштрасса} в $\overline{B}_R(0)$ у $f$ есть точка минимума на шаре $x^{*}$.
  Тогда  $\| y \| \le R \Rightarrow f(y) \ge f(x^{*}) \Rightarrow f(0) \ge f(x^{*})$.

  Тогда получается, что внутри шара это есть минимум, а вне шара $f(y) > f(0) \ge f(x^{*})$.
  Тогда $\forall y \in \R^n\ f(y) \ge f(x^{*})$.
  Тогда $x^{*}$~--- есть глобальный минимум

  \proofpart{Единственность}
  Пусть есть две точки минимума $x_1 \ne x_2$, $f(x_1) = f(x_2) =: m$.
  По определению \hyperref[note:22E546FC-186F-4A67-85AA-2C85DCEF9021]{сильной выпуклости}, $\phi(x) = f(x) - \frac{\mu}{2} \| x \|^2$  -- \hyperref[note:1CA14F2C-637F-41DE-9E97-72075D19E97F]{выпуклая}.
  По выпуклости $\phi$, $\phi(\lambda x_1 + (1-\lambda)x_2) \le \lambda \phi(x_1) + (1-\lambda) \phi(x_2)$, давайте возьмем $\lambda=\frac{1}{2}$ \\
  И пусть $z := \frac{x_1}{2} + \frac{x_2}{2}$, тогда
  \[
  \phi(z) \le \frac{1}{2} (\phi(x_1) + \phi(x_2))
  \]
  Теперь можно раскрыть значение $\phi$,
  \[
  f(z) - \frac{\mu}{2} \| z \|^2 \le \frac{1}{2} \left(  f(x_1) + f(x_2) - \frac{\mu}{2} (\| x_1\|^2 + \| x_2 \|^2)  \right)
  \]
  Но мы знаем, что $f(x_1) = f(x_2)$, можно записать их сумму, как $2m$.
  Рассмотрим $\|z\|^2$
  \begin{multline*}
  \| z\|^2 = \left\langle z; z \right\rangle = \left\langle \frac{1}{2} (x_1 + x_2); \frac{1}{2} (x_1 + x_2) \right\rangle = \\
  \frac{1}{4} \left\langle x_1 + x_2; x_1 + x_2 \right\rangle =
  \frac{1}{4} \left( \| x_1 \|^2 + 2 \left\langle x_1; x_2 \right\rangle +\| x_2 \|^2 \right)
  \end{multline*}
  Давайте подставим все наши преобразования в неравенство
  \[
  f(z) - \frac{\mu}{8} (\| x_1 \|^2 + 2 \left\langle x_1; x_2 \right\rangle +\| x_2 \|^2) \le m - \frac{\mu}{4} (\|x_1\|^2 + \|x_2\|^2)
  \]
  Давайте теперь сгруппируем члены с $\mu$
  \[
  f(z) \le m - \frac{\mu}{8} (2 \| x_1\|^2 + 2\|x_2\|^2 - \|x_1\|^2 - 2 \left\langle x_1;x_2 \right\rangle - \| x_2 \|^2) = m - \frac{\mu}{8} \| x_1 - x_2 \|^{2}
  \]

  При этом $\mu > 0$, а еще $x_1 \ne x_2$, следовательно их разность нулевой быть не может, тогда $\| x_1 - x_2 \|^2 > 0$ \\
  Из этого получаем, что
  \[
  f(z) \le m - \frac{\mu}{8} \| x_1 - x_2 \|^{2} < m
  \]
  Получили противоречие тому, что $m$~--- минимальное значение.
\end{proof}
